\documentclass[12pt,a4paper]{ctexart}

% ============================================================
%                       宏包设置
% ============================================================

\usepackage{amsmath, amssymb, amsfonts}
\usepackage{geometry}
\usepackage{fancyhdr}
\usepackage{graphicx}
\usepackage{xcolor}
\usepackage{setspace}
\usepackage[most]{tcolorbox}
\usepackage{enumitem}

% ============================================================
%                 清新蓝绿色配色
% ============================================================

% 主色：清爽青绿色
\definecolor{SZURed}{HTML}{2F9E8F}

% 题目区域：浅薄荷绿
\definecolor{SZULightRed}{HTML}{EAF7F4}

% 答案区域：浅天空蓝
\definecolor{SZULightRose}{HTML}{EEF7FC}

% 边框：柔和蓝绿色
\definecolor{SZUSoftBorder}{HTML}{9CCFC8}

% 强调文字：沉稳蓝绿色
\definecolor{SZUTextRed}{HTML}{256B72}

% ============================================================
%                       页面设置
% ============================================================

\geometry{
    left=2.5cm,
    right=2.5cm,
    top=2.5cm,
    bottom=2.5cm
}

\setlength{\parindent}{0pt}
\setlength{\parskip}{0.6em}
\linespread{1.15}

% ============================================================
%                       课程信息
% ============================================================

\newcommand{\coursename}{概率论与数理统计}
\newcommand{\department}{深圳大学生物医学工程}
\newcommand{\homeworknumber}{作业 4}

% ============================================================
%                  学生信息 —— 请学生修改
% ============================================================

\newcommand{\studentname}{请填写姓名}
\newcommand{\studentid}{请填写学号}

% ============================================================
%                       页眉页脚
% ============================================================

\pagestyle{fancy}
\fancyhf{}

\lhead{\textcolor{SZUTextRed}{\coursename}}
\rhead{\textcolor{SZUTextRed}{\homeworknumber}}
\cfoot{\textcolor{SZUTextRed}{\thepage}}

\renewcommand{\headrulewidth}{0.4pt}
\renewcommand{\headrule}{
    \hbox to\headwidth{
        \color{SZUSoftBorder}\leaders\hrule height \headrulewidth\hfill
    }
}

\setlength{\headheight}{15pt}

% ============================================================
%                    题目编号与环境设置
% ============================================================

\newcounter{problem}

% -------------------------
% 题目环境
% -------------------------

\newenvironment{problem}[1][]
{
    \stepcounter{problem}
    \begin{tcolorbox}[
        title={题目 \theproblem\if&#1&\relax\else\quad[#1]\fi},
        colback=SZULightRed,
        colframe=SZURed,
        coltitle=white,
        fonttitle=\bfseries,
        boxrule=0.8pt,
        arc=2mm,
        left=3mm,
        right=3mm,
        top=2.5mm,
        bottom=2.5mm,
        breakable
    ]
}
{
    \end{tcolorbox}
}

% -------------------------
% 答案环境
% -------------------------

\newenvironment{answer}
{
    \begin{tcolorbox}[
        title={学生答案},
        colback=SZULightRose,
        colframe=SZUSoftBorder,
        coltitle=SZUTextRed,
        fonttitle=\bfseries,
        boxrule=0.7pt,
        arc=2mm,
        left=3mm,
        right=3mm,
        top=2.5mm,
        bottom=2.5mm,
        breakable
    ]
}
{
    \end{tcolorbox}
}

% ============================================================
%                         正文
% ============================================================

\begin{document}

% ============================================================
%                         标题
% ============================================================

\begin{center}

    {\LARGE\bfseries\textcolor{SZURed}{\coursename}}

    \vspace{0.4em}

    {\large \department}

    \vspace{1em}

    {\Large\bfseries\textcolor{SZUTextRed}{\homeworknumber}}

\end{center}

\vspace{1em}

% ============================================================
%                       学生信息
% ============================================================

\begin{center}

\begin{tabular}{rl@{\hspace{3cm}}rl}

\textcolor{SZUTextRed}{\bfseries 姓名：}
&
\studentname
&
\textcolor{SZUTextRed}{\bfseries 学号：}
&
\studentid

\end{tabular}

\end{center}

\vspace{0.7em}

{\color{SZUSoftBorder}\hrule}

\vspace{1.5em}



% ============================================================
% ============================================================
%
%                           题目 1
%
% ============================================================
% ============================================================

\begin{problem}[课本习题3]

有3只球，4个盒子，盒子的编号为1,2,3,4。将球逐个独立地，随机地放入4个盒子中去。以$X$表示其中至少有一只球的盒子的最小号码（例如$X=3$表示第1号，第2号盒子是空的，第3个盒子至少有一只球），试求$E(X)$。
\end{problem}


% ============================================================
%
%              ↓↓↓↓↓↓ 学生答案区域 ↓↓↓↓↓↓
%
%          请只在 answer 环境中填写你的答案
%
% ============================================================

\begin{answer}

% ==========================================================
%                 请从下一行开始填写答案
% ==========================================================






% ==========================================================
%                 请在上一行结束填写答案
% ==========================================================

\end{answer}

% ============================================================
%
%              ↑↑↑↑↑↑ 学生答案区域 ↑↑↑↑↑↑
%
% ============================================================


\vspace{1em}



% ============================================================
% ============================================================
%
%                           题目 2
%
% ============================================================
% ============================================================

\begin{problem}[课本习题7]

\begin{enumerate}[label=(\arabic*)]
    \item 设随机变量$X$的概率密度为
\[
f(x)=
\begin{cases}
\mathrm{e}^{-x}, & x>0,\\
0, & x\le 0.
\end{cases}
\]
求(i)$Y=2X$；(ii)$Y=\mathrm{e}^{-2X}$的数学期望。
    \item 设随机变量$X_1,X_2,\dots,X_n$相互独立，且都服从$(0,1)$上的均匀分布 (i)求$U=\max\{X_1,X_2,\dots,X_n\}$的数学期望，(ii)求$V=\min\{X_1,X_2,\dots,X_n\}$的数学期望。
\end{enumerate}
\end{problem}


% ============================================================
%
%              ↓↓↓↓↓↓ 学生答案区域 ↓↓↓↓↓↓
%
% ============================================================

\begin{answer}

% ==========================================================
%                 请从下一行开始填写答案
% ==========================================================





% ==========================================================
%                 请在上一行结束填写答案
% ==========================================================

\end{answer}

% ============================================================
%
%              ↑↑↑↑↑↑ 学生答案区域 ↑↑↑↑↑↑
%
% ============================================================


\vspace{1em}



% ============================================================
% ============================================================
%
%                           题目 3
%
% ============================================================
% ============================================================

\begin{problem}[课本习题8]

设随机变量$(X,Y)$的分布律为
\[
\begin{array}{c|ccc}
Y\backslash X & 1 & 2 & 3 \\
\hline
-1 & 0.2 & 0.1 & 0.0 \\
0 & 0.1 & 0.0 & 0.3 \\
1 & 0.1 & 0.1 & 0.1
\end{array}
\]
\begin{enumerate}[label=(\arabic*)]
    \item 求$E(X),E(Y)$。
    \item 设$Z=Y/X$，求$E(Z)$。
    \item 设$Z=(X-Y)^2$，求$E(Z)$。
\end{enumerate}
\end{problem}


% ============================================================
%
%              ↓↓↓↓↓↓ 学生答案区域 ↓↓↓↓↓↓
%
% ============================================================

\begin{answer}

% ==========================================================
%                 请从下一行开始填写答案
% ==========================================================





% ==========================================================
%                 请在上一行结束填写答案
% ==========================================================

\end{answer}

% ============================================================
%
%              ↑↑↑↑↑↑ 学生答案区域 ↑↑↑↑↑↑
%
% ============================================================


\vspace{1em}



% ============================================================
% ============================================================
%
%                           题目 4
%
% ============================================================
% ============================================================

\begin{problem}[课本习题9]

\begin{enumerate}[label=(\arabic*)]
    \item 设随机变量$(X,Y)$的概率密度为
\[
f(x,y)=
\begin{cases}
12y^2, & 0\le y\le x\le 1,\\
0, & \text{其他}.
\end{cases}
\]
求$E(X),E(Y),E(XY),E(X^2+Y^2)$。
    \item 设随机变量$X,Y$的联合密度为
\[
f(x,y)=
\begin{cases}
\dfrac1y \mathrm{e}^{-(y+x/y)}, & x>0,y>0,\\
0, & \text{其他}.
\end{cases}
\]
求$E(X),E(Y),E(XY)$。
\end{enumerate}
\end{problem}


% ============================================================
%
%              ↓↓↓↓↓↓ 学生答案区域 ↓↓↓↓↓↓
%
% ============================================================

\begin{answer}

% ==========================================================
%                 请从下一行开始填写答案
% ==========================================================





% ==========================================================
%                 请在上一行结束填写答案
% ==========================================================

\end{answer}

% ============================================================
%
%              ↑↑↑↑↑↑ 学生答案区域 ↑↑↑↑↑↑
%
% ============================================================


\vspace{1em}



% ============================================================
% ============================================================
%
%                           题目 5
%
% ============================================================
% ============================================================

\begin{problem}[课本习题22(2)]

设随机变量$X,Y$相互独立，且$X\sim N(720,30^2)$，$Y\sim N(640,25^2)$，求$Z_1=2X+Y$，$Z_2=X-Y$的分布，并求概率$P\{X>Y\}$，$P\{X+Y>1400\}$。
\end{problem}


% ============================================================
%
%              ↓↓↓↓↓↓ 学生答案区域 ↓↓↓↓↓↓
%
% ============================================================

\begin{answer}

% ==========================================================
%                 请从下一行开始填写答案
% ==========================================================





% ==========================================================
%                 请在上一行结束填写答案
% ==========================================================

\end{answer}

% ============================================================
%
%              ↑↑↑↑↑↑ 学生答案区域 ↑↑↑↑↑↑
%
% ============================================================


\vspace{1em}



% ============================================================
% ============================================================
%
%                           题目 6
%
% ============================================================
% ============================================================

\begin{problem}[课本习题27]

下列各对随机变量$X$和$Y$，问哪几对是相互独立的？哪几对是不相关的。
\begin{enumerate}[label=(\arabic*)]
    \item $X\sim U(0,1),Y=X^2$.
    \item $X\sim U(-1,1),Y=X^2$.
    \item $X=\cos V,Y=\sin V,V\sim U(0,2\pi)$.
    \item 若$(X,Y)$的概率密度为$f(x,y)$，
\[
f(x,y)=
\begin{cases}
x+y, & 0<x<1,0<y<1,\\
0, & \text{其他}.
\end{cases}
\]
    \item
\[
f(x,y)=
\begin{cases}
2y, & 0<x<1,0<y<1,\\
0, & \text{其他}.
\end{cases}
\]
\end{enumerate}
\end{problem}


% ============================================================
%
%              ↓↓↓↓↓↓ 学生答案区域 ↓↓↓↓↓↓
%
% ============================================================

\begin{answer}

% ==========================================================
%                 请从下一行开始填写答案
% ==========================================================





% ==========================================================
%                 请在上一行结束填写答案
% ==========================================================

\end{answer}

% ============================================================
%
%              ↑↑↑↑↑↑ 学生答案区域 ↑↑↑↑↑↑
%
% ============================================================


\vspace{1em}



% ============================================================
% ============================================================
%
%                           题目 7
%
% ============================================================
% ============================================================

\begin{problem}[课本习题31]

设随机变量$(X,Y)$具有概率密度
\[
f(x,y)=
\begin{cases}
1, & |y|<x,0<x<1,\\
0, & \text{其他}.
\end{cases}
\]
求$E(X),E(Y),\mathrm{Cov}(X,Y)$。
\end{problem}


% ============================================================
%
%              ↓↓↓↓↓↓ 学生答案区域 ↓↓↓↓↓↓
%
% ============================================================

\begin{answer}

% ==========================================================
%                 请从下一行开始填写答案
% ==========================================================





% ==========================================================
%                 请在上一行结束填写答案
% ==========================================================

\end{answer}

% ============================================================
%
%              ↑↑↑↑↑↑ 学生答案区域 ↑↑↑↑↑↑
%
% ============================================================


\vspace{1em}



% ============================================================
% ============================================================
%
%                           题目 8
%
% ============================================================
% ============================================================

\begin{problem}[课本习题32]

设随机变量$(X,Y)$具有概率密度
\[
f(x,y)=
\begin{cases}
\dfrac18(x+y), & 0\le x\le 2,0\le y\le 2,\\
0, & \text{其他}.
\end{cases}
\]
求$E(X),E(Y),\mathrm{Cov}(X,Y),\rho_{XY},D(X+Y)$。
\end{problem}


% ============================================================
%
%              ↓↓↓↓↓↓ 学生答案区域 ↓↓↓↓↓↓
%
% ============================================================

\begin{answer}

% ==========================================================
%                 请从下一行开始填写答案
% ==========================================================





% ==========================================================
%                 请在上一行结束填写答案
% ==========================================================

\end{answer}

% ============================================================
%
%              ↑↑↑↑↑↑ 学生答案区域 ↑↑↑↑↑↑
%
% ============================================================





\end{document}
```
